% Irodov, Arihant PDF pp. 13--15; original printed pages 4--6.
% Chinese translation and independently derived worked solutions: problems 21--30.
\exproblem{21}{4}
一质点在 $t=0$ 时从原点沿 $x$ 轴正方向出发，其速度随时间变化为
$v_x=v_0(1-t/\tau)$，其中 $v_0=10.0\,\mathrm{cm/s}$，$\tau=5.0\,\mathrm s$。求：
(a) $t=6.0,10,20\,\mathrm s$ 时的坐标；
(b) 质点距原点 $10.0\,\mathrm{cm}$ 的所有非负时刻；
(c) 前 $4.0\,\mathrm s$ 和 $8.0\,\mathrm s$ 的路程，并作 $s(t)$ 的定性图。
\exsolution
直接积分并利用 $x(0)=0$ 得
\[
 x(t)=\int_0^t v_0(1-u/\tau)\,\mathrm du
 =v_0t\Bigl(1-\frac{t}{2\tau}\Bigr).
\]
以 $\mathrm{cm}$ 和 $\mathrm s$ 计，$x(t)=10t-t^2$，故
\[
 \boxed{x(6)=24\,\mathrm{cm},\quad x(10)=0,\quad
 x(20)=-200\,\mathrm{cm}=-2.00\,\mathrm m.}
\]
距原点 $10\,\mathrm{cm}$ 意味着 $|x(t)|=10\,\mathrm{cm}$，而非只取 $x=10$。
\[
 t^2-10t+10=0\quad\text{或}\quad t^2-10t-10=0.
\]
三组非负根为
\[
 \boxed{t=5-\sqrt{15}\simeq1.13\,\mathrm s,\ \
 5+\sqrt{15}\simeq8.87\,\mathrm s,\ \
 5+\sqrt{35}\simeq10.92\,\mathrm s.}
\]
由于 $t=\tau=5\,\mathrm s$ 时速度反号，不能用位移代替全程路程：
\[
 s(t)=\begin{cases}
 v_0t(1-t/(2\tau)),&0\le t\le\tau,\\
 v_0\tau-v_0t(1-t/(2\tau)),&t\ge\tau .
 \end{cases}
\]
代入得 $\boxed{s(4)=24\,\mathrm{cm},\ s(8)=34\,\mathrm{cm}}$。
曲线前段凹向下，在 $t=5\,\mathrm s$ 处斜率为零，此后凹向上、路程继续增加。
\excheck 对照本版书末答案，第21(a)题的末项印为 $-4.0\,\mathrm m$，与给定
$v_0,\tau,t$ 不相容；依题面逐项代入应为 $-2.00\,\mathrm m$，记录为版本勘误。

\exproblem{22}{4--5}
质点沿 $x$ 轴正方向运动，速率满足 $v=\alpha\sqrt{x}$，$\alpha>0$，且
$x(0)=0$。求：(a) 速率及加速度的时间规律；(b) 质点走过最初 $s$ 米期间的平均速率。
\exsolution
对离开原点后 $x>0$ 的运动分离变量：
\[
 \frac{\mathrm dx}{\sqrt{x}}=\alpha\,\mathrm dt,\qquad
 2\sqrt{x}=\alpha t,
\quad x=\frac{\alpha^2t^2}{4}.
\]
因此
\[
 \boxed{v(t)=\frac{\alpha^2t}{2},\qquad a_x(t)=\frac{\alpha^2}{2}.}
\]
走完路程 $s$ 所需时间为 $T=2\sqrt{s}/\alpha$，所以
\[
 \boxed{\bar v=\frac{s}{T}=\frac{\alpha\sqrt{s}}2.}
\]
\excheck 本题应补充“$t=0$ 后立即开始移动”的物理选解。
数学上 $\dot x=\alpha\sqrt{x}$ 在 $x=0$ 不满足局部 Lipschitz 条件，
还允许先静止任意时间 $t_*$，再取
$x(t)=\alpha^2(t-t_*)^2/4$；单独给出 $x(0)=0$ 不保证唯一性。

\exproblem{23}{5}
质点作直线减速运动，减速度的大小为 $a\sqrt v$，$a>0$ 为常数，
初速 $v_0>0$。求停止前走过的总路程以及停止所需的时间。
\exsolution
在 $v>0$ 阶段，
\[
 \frac{\mathrm dv}{\mathrm dt}=-a\sqrt v,\qquad
 2\sqrt v=2\sqrt{v_0}-at.
\]
当 $v=0$ 时
\[
 \boxed{T=\frac{2\sqrt{v_0}}a.}
\]
用 $\mathrm dt=-\mathrm dv/(a\sqrt v)$ 算路程：
\[
 s=\int_0^T v\,\mathrm dt
   =\frac1a\int_0^{v_0}\sqrt v\,\mathrm dv
   =\boxed{\frac{2v_0^{3/2}}{3a}}.
\]
\excheck $[a]=\mathrm{m}^{1/2}\mathrm{s}^{-3/2}$，
则 $[\sqrt{v_0}/a]=\mathrm s$，$[v_0^{3/2}/a]=\mathrm m$。

\exproblem{24}{5}
质点的位矢为 $\bm r(t)=at\,\bm i-bt^2\,\bm j$，
其中 $a,b>0$。求：(a) 轨迹 $y(x)$；(b) 速度、加速度的矢量及大小；
(c) 二者夹角 $\theta(t)$；(d) 前 $t$ 秒的平均速度矢量及大小。
\exsolution
消去 $t=x/a$，有
\[
 \boxed{y=-\frac b{a^2}x^2},
\]
即过原点且开口朝下的抛物线。
对位矢求导：
\[
 \boxed{\bm v=a\bm i-2bt\bm j,\qquad
 \bm a=-2b\bm j},\quad
 |\bm v|=\sqrt{a^2+4b^2t^2},\quad|\bm a|=2b.
\]
当 $t\ge0$ 时，通过点积计算夹角：
\[
 \cos\theta=\frac{\bm v\cdot\bm a}{|\bm v||\bm a|}
 =\frac{2bt}{\sqrt{a^2+4b^2t^2}},\qquad
 \boxed{\tan\theta=\frac a{2bt}\ (t>0)}.
\]
$t=0$ 时 $\theta=\pi/2$。前 $t$ 秒平均速度为
\[
 \boxed{\bar{\bm v}(t)=\frac{\bm r(t)-\bm r(0)}t
 =a\bm i-bt\bm j,\quad
 |\bar{\bm v}|=\sqrt{a^2+b^2t^2}.}
\]
\excheck 轨迹开口、加速度方向和末态平均速度的负 $y$ 分量互相吻合。

\exproblem{25}{5}
质点运动方程为 $x=at,\ y=at(1-\alpha t)$，
$a,\alpha>0$。求：(a) 轨迹方程；(b) 速度与加速度；
(c) 速度和加速度夹角为 $\pi/4$ 的时刻。
\exsolution
用 $t=x/a$ 消去时间，得到
\[
 \boxed{y=x-\frac{\alpha}{a}x^2}.
\]
微分可得
\[
 \boxed{\bm v=a\bm i+a(1-2\alpha t)\bm j,\qquad
 \bm a=-2a\alpha\bm j}.
\]
加速度沿负 $y$ 轴，所以夹角为锐角 $\pi/4$ 需要
$v_y<0$ 且 $|v_y|=v_x=a$，即 $1-2\alpha t=-1$，于是
\[
 \boxed{t_0=\frac1\alpha}.
\]
\excheck $t_0$ 时 $\bm v=(a,-a)$，与 $\bm a=(0,-2a\alpha)$
的夹角确为 $45^\circ$；不能误取 $t=0$，那时夹角为 $135^\circ$。

\exproblem{26}{5}
质点的平面运动方程为 $x=a\sin\omega t$，
$y=a(1-\cos\omega t)$，其中 $a,\omega>0$。
求：(a) 经过时间 $\tau$ 的路程；(b) 速度与加速度的夹角。
\exsolution
因 $x^2+(y-a)^2=a^2$，运动轨迹是半径为 $a$、
圆心在 $(0,a)$ 的圆。微分：
\[
 \bm v=a\omega(\cos\omega t\,\bm i+\sin\omega t\,\bm j),\quad
 \bm a=a\omega^2(-\sin\omega t\,\bm i+\cos\omega t\,\bm j).
\]
速度大小恒为 $a\omega$，且 $\bm v\cdot\bm a=0$，
因此
\[
 \boxed{s=a\omega\tau,\qquad \angle(\bm v,\bm a)=\frac\pi2.}
\]
\excheck 周期 $2\pi/\omega$ 中路程 $2\pi a$，与圆周长一致。

\exproblem{27}{5}
质点在 $xy$ 平面内受沿负 $y$ 方向的恒定加速度，大小为 $w$，
其运动轨迹为 $y=Ax-bx^2$，$A,b>0$。求经过坐标原点时的速度大小。
\exsolution
取原点对应的时刻为 $t=0$，记初速度分量为 $v_{x0},v_{y0}$。
因为 $a_x=0$、$a_y=-w$，
\[
 x=v_{x0}t,\qquad y=v_{y0}t-\frac12wt^2
 =\frac{v_{y0}}{v_{x0}}x-\frac{w}{2v_{x0}^2}x^2.
\]
与题给轨迹逐项比较：
\[
 \frac{v_{y0}}{v_{x0}}=A,\qquad
 \frac{w}{2v_{x0}^2}=b.
\]
于是
\[
 \boxed{v_{x0}=\sqrt{\frac w{2b}},\quad
 v_{y0}=A\sqrt{\frac w{2b}},\quad
 v_0=\sqrt{\frac{w(1+A^2)}{2b}}.}
\]
\excheck $b$ 的量纲为 $\mathrm m^{-1}$，根号内量纲为
$\mathrm m^2/\mathrm s^2$。若取反向通过原点，相应分量可同时变号而速率不变。

\exproblem{28}{5}
一小物体以初速度矢量 $\bm v_0$ 与水平面成角抛出，忽略空气阻力。
以重力加速度矢量 $\bm g$ 表示：
(a) 位移 $\bm r(t)$；(b) 前 $t$ 秒的平均速度，以及回到发射高度时全程的平均速度。
\exsolution
以抛出点为位移原点，匀加速运动给出
\[
 \boxed{\bm r(t)=\bm v_0t+\frac12\bm g\,t^2,\qquad
 \bar{\bm v}(0,t)=\bm v_0+\frac12\bm g\,t.}
\]
若初速度与重力方向夹角大于 $90^\circ$，抛体能重新穿过发射高度。
由竖直分量归零得飞行时间
\[
 T=-\frac{2(\bm v_0\cdot\bm g)}{g^2}>0.
\]
故
\[
 \boxed{\bar{\bm v}_{\rm flight}=
 \bm v_0-\frac{\bm g(\bm v_0\cdot\bm g)}{g^2}},
\]
正好是初速度在垂直于 $\bm g$ 的平面内的投影。
\excheck 此式没有把沿抛物线的总路程与位移混为一谈。

\exproblem{29}{5--6}
由地面以初速 $v_0$、与水平夹角 $\alpha$ 抛出一物体，忽略空气阻力。
求：(a) 飞行时间；(b) 最大高度、水平射程，二者何时相等；
(c) 轨迹方程；(d) 发射点与最高点处的曲率半径。
\exsolution
取 $x$ 水平、$y$ 竖直向上：
\[
 x=v_0\cos\alpha\,t,\qquad
 y=v_0\sin\alpha\,t-\frac12gt^2.
\]
回到 $y=0$ 的非零根给出
\[
 \boxed{T=\frac{2v_0\sin\alpha}{g},\quad
 H=\frac{v_0^2\sin^2\alpha}{2g},\quad
 L=\frac{v_0^2\sin2\alpha}{g}.}
\]
$H=L$ 且 $0<\alpha<\pi/2$ 时有
$\frac12\sin\alpha=2\cos\alpha$，即
\[
 \boxed{\tan\alpha=4,\quad\alpha=\arctan4\simeq75.96^\circ.}
\]
消元得
\[
 \boxed{y=x\tan\alpha-
 \frac{gx^2}{2v_0^2\cos^2\alpha}}.
\]
平面曲线的曲率半径 $R=v^3/|\bm v\times\bm g|$，
因为 $|\bm v\times\bm g|=gv_0\cos\alpha$，
所以发射点速度为 $v_0$、顶点速度为 $v_0\cos\alpha$ 时分别为
\[
 \boxed{R_{\rm launch}=\frac{v_0^2}{g\cos\alpha},\qquad
 R_{\rm top}=\frac{v_0^2\cos^2\alpha}{g}.}
\]
\excheck 当 $\alpha\to0$ 时 $H\to0$ 且 $L\to0$；
$R_{\rm top}$ 有限但最高点与地面重合，属极限情况。

\exproblem{30}{6}
沿用第29题条件，画出法向加速度大小 $a_n$、切向加速度大小 $|a_\tau|$
以及总加速度在速度方向上的有符号投影 $a_v$ 随时间的大致变化。
\exsolution
速度分量 $v_x=v_0\cos\alpha$、$v_y=v_0\sin\alpha-gt$，
速率 $v(t)=\sqrt{v_x^2+v_y^2}$。用 $\bm a=(0,-g)$ 的点积和叉积，
\[
 \boxed{a_v(t)=a_\tau(t)=-\frac{g(v_0\sin\alpha-gt)}{v(t)},\quad
 |a_\tau|=\frac{g|v_0\sin\alpha-gt|}{v(t)},\quad
 a_n(t)=\frac{gv_0\cos\alpha}{v(t)}.}
\]
在 $t=0$ 时 $a_v=-g\sin\alpha,\ a_n=g\cos\alpha$；
在顶点 $t=T/2$ 时 $a_v=0,\ a_n=g$；
落地前 $a_v=g\sin\alpha,\ a_n=g\cos\alpha$。
故 $a_n(t)$ 关于 $T/2$ 对称、先增后减；
$|a_\tau|$ 先减至零再增；
$a_v$ 从负值连续增加、过零后为正。
\excheck 任意时刻 $a_\tau^2+a_n^2=g^2$，图形的纵坐标不应超过 $g$。
